\documentclass[10pt]{report}
\usepackage{anyfontsize}
\AtBeginDocument{\fontsize{8.5}{10.6}\selectfont}

% ===================== ENGINE GUARD =====================
\usepackage{iftex}
\ifXeTeX\else
  \errmessage{This document must be compiled with XeLaTeX.}
\fi

% ===================== PAGE LAYOUT =====================
\usepackage[top=2cm,bottom=2.5cm,left=3cm,right=2cm]{geometry}
\setlength{\headheight}{15pt}

% ===================== CORE PACKAGES =====================
\usepackage{graphicx}
\usepackage{float}
\usepackage{array}
\usepackage{tabularx}
\usepackage{booktabs}
\usepackage{multirow}
\usepackage{threeparttable}
\usepackage{colortbl}
\usepackage{enumitem}
\usepackage{amsmath}
\usepackage{amssymb}
\usepackage{ragged2e}

% ===================== VISUAL DESIGN =====================
\usepackage[table]{xcolor}
\usepackage{tcolorbox}
\tcbuselibrary{breakable,skins,theorems}
\usepackage{framed}

\usepackage{fancyhdr}
\usepackage{titlesec}
\usepackage{tocloft}

% ===================== HEADER/FOOTER STYLING =====================
\pagestyle{fancy}
\fancyhf{}
\fancyhead[RE]{\nouppercase{\rightmark}}
\fancyhead[LO]{\nouppercase{\leftmark}}
\fancyhead[LE,RO]{\thepage}
\renewcommand{\headrulewidth}{0.4pt}
\renewcommand{\footrulewidth}{0pt}
\fancypagestyle{plain}{%
  \fancyhf{}
  \fancyfoot[C]{\thepage}
  \renewcommand{\headrulewidth}{0pt}
}

% ===================== TIKZ =====================
\usepackage{tikz}
\usetikzlibrary{shapes,arrows,positioning,calc,decorations.pathreplacing}

% ===================== PERSIAN (XeLaTeX) =====================
\usepackage{xepersian}

% ============ LATIN FONT SETUP ============
\setlatintextfont[
  Scale=0.95,
  Numbers=Lining,
  Ligatures=TeX
]{TeX Gyre Termes}

% ============ PERSIAN FONT SETUP ============
\settextfont[
  Path=fonts/,
  Extension=.ttf,
  UprightFont=*,
  BoldFont=*Bd,
  ItalicFont=*It,
  BoldItalicFont=*BdIt,
  Script=Arabic,
  Language=Persian
]{XB-Yagut}

\setdigitfont[
  Path=fonts/,
  Extension=.ttf,
  UprightFont=*,
  BoldFont=*Bd,
  ItalicFont=*It,
  BoldItalicFont=*BdIt,
  Numbers=Proportional,
  Script=Arabic,
  Language=Persian
]{XB-Yagut}

% ============ CHAPTER/TITLE FORMATTING ============
\titleformat{\chapter}[display]
  {\normalfont\huge\bfseries\color{headerColor}}
  {\flushright\chaptertitlename\ \thechapter}
  {20pt}
  {\Huge\flushright}
\titlespacing*{\chapter}{0pt}{-20pt}{30pt}

\titleformat{\section}
  {\normalfont\Large\bfseries\color{headerColor}}
  {\thesection}{1em}{}

% ============ TABLE OF CONTENTS STYLING ============
\renewcommand{\cftchapfont}{\bfseries\color{headerColor}}
\renewcommand{\cftchapleader}{\cftdotfill{\cftdotsep}}
\setlength{\cftbeforechapskip}{6pt}
\renewcommand{\contentsname}{\hfill\Large\bfseries\color{headerColor}فهرست مطالب\hfill}

% ============ ENHANCED MACROS FOR BILINGUAL CONTENT ============
\newcommand{\lat}[1]{\lr{#1}}
\newcommand{\latinterm}[1]{\lr{\textit{#1}}}
\newcommand{\latinbold}[1]{\lr{\textbf{#1}}}
\newcommand{\latinnum}[1]{\lr{#1}}
\newcommand{\technical}[1]{\lr{\textsc{#1}}}

% ============ TABLE COLUMN TYPES FOR RTL ============
\newcolumntype{P}[1]{>{\RaggedRight}p{#1}}
\newcolumntype{H}[1]{>{\bfseries\color{headerColor}\RaggedRight}p{#1}}
\newcolumntype{C}[1]{>{\centering\arraybackslash}p{#1}}

% ============ LINE SPACING ============
\linespread{1.15}

% ============ HYPERREF SETUP ============
\usepackage{hyperref}
\hypersetup{
    colorlinks=true,
    linkcolor=headerColor,
    citecolor=headerColor,
    filecolor=headerColor,
    urlcolor=headerColor,
    pdftitle={لیست تماس اعضای هیئت علمی دانشگاه صنعتی شریف},
    pdfauthor={دانشگاه صنعتی شریف},
    pdfsubject={اطلاعات تماس اساتید},
    pdfkeywords={ایمیل، تماس، هیئت علمی، شریف},
    bookmarksopen=true,
    pdfpagemode=UseOutlines
}

% ===================== PROFESSIONAL COLOR SCHEME =====================
\definecolor{headerColor}{RGB}{44,62,80}
\definecolor{accentColor}{RGB}{90, 70, 70}
\definecolor{pathwayIncome}{RGB}{235, 242, 248}
\definecolor{pathwayEquity}{RGB}{248, 243, 238}
\definecolor{pathwayMarket}{RGB}{238, 246, 240}
\definecolor{pathwayAssets}{RGB}{250, 248, 238}
\definecolor{tableHeader}{RGB}{44, 62, 80}
\definecolor{tableRow1}{RGB}{255, 255, 255}
\definecolor{tableRow2}{RGB}{248, 249, 250}

% ===================== PROFESSIONAL TCOLORBOX STYLES =====================
\newtcolorbox{keyconceptbox}[1][]{
  enhanced,
  colback=pathwayIncome,
  colframe=headerColor,
  fonttitle=\bfseries,
  title={\color{white}مفهوم کلیدی},
  coltitle=white,
  attach boxed title to top right={yshift=-2mm, xshift=-4mm},
  boxed title style={colback=headerColor, sharp corners, size=small},
  breakable,
  sharp corners,
  boxrule=0.75pt,
  left=4mm,
  right=4mm,
  top=3mm,
  bottom=3mm,
  #1
}

\newtcolorbox{analysisbox}[2][]{
  enhanced,
  colback=pathwayMarket,
  colframe=headerColor,
  fonttitle=\bfseries,
  title={#2},
  breakable,
  sharp corners,
  boxrule=0.75pt,
  left=4mm,
  right=4mm,
  top=3mm,
  bottom=3mm,
  #1
}

\newtcolorbox{infobox}[2][]{
  enhanced,
  colback=pathwayAssets,
  colframe=headerColor,
  fonttitle=\bfseries,
  title={#2},
  breakable,
  sharp corners,
  boxrule=0.75pt,
  left=4mm,
  right=4mm,
  top=3mm,
  bottom=3mm,
  #1
}

% ===================== DOCUMENT STARTS =====================
\begin{document}

% ===================== TITLE PAGE =====================
\thispagestyle{empty}
\begin{center}
\vspace*{2cm}

{\Huge\bfseries\color{headerColor} اطلاعات تماس اعضای هیئت علمی و تحلیل استراتژیک}

\vspace{0.6cm}
{\color{gray}\rule{0.7\textwidth}{0.75pt}}
\vspace{0.6cm}

{\Large\bfseries دانشگاه صنعتی شریف}

\vspace{0.2cm}
{\large دانشکده‌های مهندسی برق و فیزیک}

\vspace{1.5cm}

\begin{tikzpicture}[scale=1.2, transform shape]
    \node[draw=headerColor, fill=white, rectangle, rounded corners, minimum width=6cm, minimum height=1.5cm, line width=1.5pt] 
    at (0,0) {\Large\bfseries \lat{Sharif University}};
    
    \node[draw=headerColor, fill=white, rectangle, rounded corners, minimum width=5cm, minimum height=1cm, line width=1pt, yshift=-1.5cm] 
    at (0,0) {\large\bfseries \lat{Faculty Contact Directory \& Analysis}};
\end{tikzpicture}

\vfill
{\small\color{gray} تاریخ انتشار: ۱۰ بهمن ۱۴۰۴}
\end{center}

% ===================== TABLE OF CONTENTS =====================
\clearpage
\tableofcontents

% ===================== CHAPTER 1: STRATEGIC ANALYSIS =====================
\clearpage
\chapter{تحلیل استراتژیک: دانشگاه صنعتی شریف و ژئوپلیتیک نیمه‌هادی‌ها}

\begin{keyconceptbox}
این بخش شامل تحلیل جامع دانشگاه صنعتی شریف، نقش آن در ژئوپلیتیک نیمه‌هادی‌ها و همکاری‌های بین‌المللی بر اساس اطلاعات ارائه شده می‌باشد.
\end{keyconceptbox}

\section{پیشینه دانشگاه صنعتی شریف (SUT) و برتری آن در مهندسی برق}

دانشگاه صنعتی شریف برجسته‌ترین مؤسسه فنی ایران است. دانشکده مهندسی برق (EE) آن دارای \textbf{بیش از ۷۰ عضو هیئت علمی} است و حدود \textbf{۸۵۰ دانشجوی کارشناسی و ۶۵۰ دانشجوی تحصیلات تکمیلی} را آموزش می‌دهد که آن را به یکی از بزرگ‌ترین دانشکده‌های مهندسی در ایران تبدیل کرده است. پذیرش در آن بسیار انتخابی است — دانشجویان رتبه‌برتر از آزمون سراسری و مدال‌آوران المپیاد به‌طور معمول SUT را انتخاب می‌کنند. این دانشکده برنامه‌های پیشرفته‌ای در مهندسی قدرت، مخابرات، مایکروویو و فوتونیک، سیستم‌های کنترل، الکترونیک، سیستم‌های دیجیتال و مهندسی پزشکی ارائه می‌دهد. این محیط آکادمیک سخت‌گیرانه، فارغ‌التحصیلانی تولید می‌کند که اغلب به تحصیلات تکمیلی در دانشگاه‌های برتر خارج از کشور می‌پردازند یا وارد موقعیت‌های صنعتی و تحقیقاتی می‌شوند.

\section{مشارکت SUT در ژئوپلیتیک نیمه‌هادی‌ها}

\subsection{برگزاری نشست‌ها در مورد ژئوپلیتیک نیمه‌هادی‌ها}

در ماه می ۲۰۲۴، مؤسسه تحقیقات سیاست علمی دانشگاه صنعتی شریف (SPRI) رویدادی سیاستی با عنوان \textbf{"ژئوپلیتیک نیمه‌هادی‌ها"} برگزار کرد. این نشست به بررسی \textit{"رقابت‌ها و تنش‌های عمده"} در صنعت جهانی نیمه‌هادی‌ها پرداخت و پیامدهای ژئوپلیتیک رقابت و همکاری میان ملت‌ها در نوآوری و تولید را بررسی کرد.

شرکت‌کنندگانی از هند، تایوان، هنگ‌کنگ و ایران در مورد چگونگی شکل‌دهی عوامل جغرافیایی — مانند دسترسی به منابع، تخصص نیروی کار و زیرساخت — به قابلیت‌های ملی در طراحی و ساخت تراشه بحث کردند. این رویداد نشان‌دهنده تلاش SUT برای گردهم آوردن متخصصان برای درک پویایی‌های بین‌المللی زنجیره‌های تأمین نیمه‌هادی‌ها و انتخاب‌های سیاستی پیش روی ملت‌ها است.

\subsection{مشارکت فارغ‌التحصیلان SUT در تحقیقات نیمه‌هادی‌ها}

فارغ‌التحصیلان SUT نقش فعالی در تحقیقات جهانی نیمه‌هادی‌ها ایفا می‌کنند. برای مثال، \textbf{پروفسور مهدی طاهری}، دانش‌آموخته کارشناسی مهندسی برق SUT، ریاست کرسی محاسبات نانومقیاس قابل اطمینان در موسسه فناوری کارلسروهه و استاد مهمان در \textbf{imec}، قطب تحقیقات ریزالکترونیک اروپا است.

طاهری در مصاحبه‌ای با imec توضیح داد که دانشگاه‌ها به \textit{"کیت‌های طراحی فرآیند (PDK) اکتشافی"} نیاز دارند که به محققان اجازه می‌دهد گره‌های ساخت جدید را قبل از در دسترس بودن PDK های تجاری بررسی کنند. چنین کیت‌هایی تحقیقات آکادمیک و کاربرد صنعتی در گره‌های پیشرفته نیمه‌هادی را به هم پیوند می‌دهند و به دانشگاه‌ها امکان می‌دهند در طراحی تراشه مرتبط باقی بمانند و با صنعت همکاری کنند.

\section{صنعت جهانی نیمه‌هادی‌ها و چالش‌های زنجیره تأمین}

صنعت نیمه‌هادی‌ها پایه فناوری مدرن است، از هوش مصنوعی و محاسبات高性能 گرفته تا خودروها و شبکه‌های ارتباطی. \textbf{پیش‌بینی می‌شود درآمد جهانی تا اواخر دهه ۲۰۲۰ از ۱ تریلیون دلار آمریکا فراتر رود}، با سرمایه‌گذاری‌های برنامه‌ریزی شده حدود ۱ تریلیون دلار برای احداث کارخانه‌های ساخت جدید تا سال ۲۰۳۰.

رشد توسط تقاضا برای شتاب‌دهنده‌های هوش مصنوعی، حافظه‌های پیشرفته و تراشه‌های سه‌بعدی انباشته شده هدایت می‌شود، اما زنجیره تأمین همچنان بسیار متمرکز و آسیب‌پذیر باقی مانده است: بیشتر تولیدات پیشرفته در شرق آسیا واقع شده‌اند، مواد اولیه حیاتی مانند ویفرهای سیلیکونی و مواد شیمیایی توسط تعداد معدودی از کشورها تأمین می‌شوند، و کمبود نیروی کار و هزینه‌های بالا، احداث کارخانه‌های ساخت (fab) خارج از آسیا را پرهزینه می‌کند.

دولت‌ها با یارانه‌ها و مشارکت‌های استراتژیک پاسخ می‌دهند تا تولید را بومی‌سازی کرده و دسترسی به تراشه‌ها و مواد را تضمین کنند.

\section{همکاری‌های بین‌المللی با مشارکت کره جنوبی}

\subsection{مشارکت دیجیتال اتحادیه اروپا-کره جنوبی و پروژه‌های مشترک}

\textbf{مشارکت دیجیتال اتحادیه اروپا–کره جنوبی} که در \textbf{نوامبر ۲۰۲۲} راه‌اندازی شد، با هدف تقویت همکاری در زمینه نیمه‌هادی‌ها، 5G/6G، محاسبات کوانتومی و هوش مصنوعی است. در \textbf{اوایل سال ۲۰۲۴}، \textbf{ابتکار مشترک تراشه‌های اتحادیه اروپا (Chips JU)} و دولت کره جنوبی فراخوانی برای پروژه‌های تحقیقاتی پیش از تجاری‌سازی صادر کردند.

چهار پروژه — \textbf{ENERGIZE}, \textbf{NEHIL}, \textbf{HAETAE} و \textbf{ViTFOX} — انتخاب شدند. هر کنسرسیوم، تیم‌های تحقیقاتی اروپایی و کره‌ای را با هم جفت می‌کند و بر یکپارچه‌سازی ناهمگن و تراشه‌های نورومورفیک متمرکز است. این پروژه‌ها \textbf{۱۲ میلیون یورو} بودجه دریافت می‌کنند، نیمی از اتحادیه اروپا و نیمی از کره جنوبی، و تسهیلات را به اشتراک گذاشته و انجمن‌های محققان نیمه‌هادی را سازماندهی خواهند کرد.

\subsection{همکاری هند–کره جنوبی و همکاری سه‌جانبه با ژاپن}

\textbf{سیاست جدید جنوبی} کره جنوبی همکاری با هند را تشویق می‌کند. نظرات حاکی از آن است که همکاری می‌تواند شامل مراکز نوآوری مشترک برای تراشه‌های هوش مصنوعی و بسته‌بندی سه‌بعدی باشد که از استعداد طراحی نیمه‌هادی هند و تخصص تولید کره جنوبی بهره می‌برد.

تحلیلگران استدلال می‌کنند که \textbf{برنامه Semicon هند} و \textbf{قانون ویژه پرورش فناوری استراتژیک ملی} و \textbf{کمربند نیمه‌هادی K} کره جنوبی می‌توانند مکمل یکدیگر باشند و زنجیره تأمین خودکفاتری را ترویج دهند. نویسندگان همچنین همکاری سه‌جانبه‌ای بین \textbf{هند، ژاپن و کره جنوبی} را پیشنهاد می‌دهند تا نقاط قوت آن‌ها (تولید در ژاپن و کره جنوبی، استعداد طراحی در هند) ادغام شود.

\section{تقاطع SUT و همکاری‌های جهانی نیمه‌هادی}

اگرچه هیچ شواهد عمومی از همکاری‌های تحقیقاتی مستقیم \textbf{دانشگاه صنعتی شریف–کره جنوبی} در زمینه نیمه‌هادی‌ها وجود ندارد، SUT از طریق تحقیق، آموزش و گفت‌وگوی سیاستی به‌طور فعال با مسائل جهانی نیمه‌هادی درگیر می‌شود:

\begin{itemize}[rightmargin=1.5cm]
\item \textbf{رهبری فکری:} با برگزاری نشست‌هایی مانند رویداد \textbf{ژئوپلیتیک نیمه‌هادی‌ها}
\item \textbf{مشارکت‌های تحقیقاتی از طریق فارغ‌التحصیلان:} فارغ‌التحصیلان SUT مانند \textbf{پروفسور مهدی طاهری}
\item \textbf{یادگیری زمینه‌ای:} همکاری‌های بین‌المللی خلاصه شده در بالا — پروژه‌های مشترک اتحادیه اروپا-کره جنوبی، همکاری هلند–کره جنوبی، انجمن‌های بریتانیا–کره جنوبی، ابتکارات هند–کره جنوبی و چارچوب‌های چندجانبه
\end{itemize}

% ===================== CHAPTER 2: ELECTRICAL ENGINEERING =====================
\clearpage
\chapter{اعضای هیئت علمی مهندسی برق}

\begin{keyconceptbox}
لیست کامل ایمیل‌های اعضای هیئت علمی دانشکده مهندسی برق دانشگاه صنعتی شریف. این اطلاعات از صفحه رسمی دانشکده استخراج شده‌اند.
\end{keyconceptbox}

\begin{table}[H]
\centering
\caption{لیست ایمیل‌های اعضای هیئت علمی مهندسی برق}
\label{tab:ee_emails}
\begin{threeparttable}
\begin{tabularx}{\textwidth}{H{4cm} X C{3.5cm}}
\toprule
\rowcolor{tableHeader}
\textcolor{white}{\textbf{نام}} & \textcolor{white}{\textbf{سمت}} & \textcolor{white}{\textbf{آدرس ایمیل}} \\
\midrule
\rowcolor{tableRow1}
دکتر فرید آشیانی & دانشیار & \href{mailto:ashtiani@sharif.edu}{\lat{ashtiani@sharif.edu}} \\
\rowcolor{tableRow2}
دکتر بهزاد آهی & استادیار & \href{mailto:ahi@sharif.edu}{\lat{ahi@sharif.edu}} \\
\rowcolor{tableRow1}
دکتر مهدی احمدی بروجنی & استادیار & \href{mailto:ahmadi@sharif.edu}{\lat{ahmadi@sharif.edu}} \\
\rowcolor{tableRow2}
دکتر فاطمه اکبر & استادیار & \href{mailto:fakbar@sharif.edu}{\lat{fakbar@sharif.edu}} \\
\rowcolor{tableRow1}
دکتر محمود اکبری & دانشیار & \href{mailto:makbari@sharif.edu}{\lat{makbari@sharif.edu}} \\
\rowcolor{tableRow2}
دکتر روح‌الله امیری & استادیار & \href{mailto:amirl@sharif.edu}{\lat{amirl@sharif.edu}} \\
\rowcolor{tableRow1}
دکتر آرش امینی & دانشیار & \href{mailto:aamini@sharif.edu}{\lat{aamini@sharif.edu}} \\
\rowcolor{tableRow2}
دکتر هاشم عرائی & استاد & \href{mailto:oraee@sharif.edu}{\lat{oraee@sharif.edu}} \\
\rowcolor{tableRow1}
دکتر مریم بابازاده & دانشیار & \href{mailto:babazadeh@sharif.ir}{\lat{babazadeh@sharif.ir}} \\
\rowcolor{tableRow2}
دکتر مسعود بابایی‌زاده & استاد & \href{mailto:mbzadeh@sharif.edu}{\lat{mbzadeh@sharif.edu}} \\
\rowcolor{tableRow1}
دکتر سعید باقری شورکی & استاد & \href{mailto:bagheri-s@sharif.edu}{\lat{bagheri-s@sharif.edu}} \\
\rowcolor{tableRow2}
دکتر علی بنایی & استاد & \href{mailto:banai@sharif.edu}{\lat{banai@sharif.edu}} \\
\rowcolor{tableRow1}
دکتر حمید بهروزی & دانشیار & \href{mailto:behroozi@sharif.edu}{\lat{behroozi@sharif.edu}} \\
\rowcolor{tableRow2}
دکتر فریدون بهنیا & استاد & \href{mailto:behnia@sharif.edu}{\lat{behnia@sharif.edu}} \\
\rowcolor{tableRow1}
دکتر محمدرضا پاکروان & دانشیار & \href{mailto:pakravan@sharif.edu}{\lat{pakravan@sharif.edu}} \\
\rowcolor{tableRow2}
دکتر مصطفی پرنیانی & استاد & \href{mailto:parniani@sharif.edu}{\lat{parniani@sharif.edu}} \\
\rowcolor{tableRow1}
دکتر محمدرضا تواژه‌ای & استاد & \href{mailto:tavazoei@sharif.edu}{\lat{tavazoei@sharif.edu}} \\
\rowcolor{tableRow2}
دکتر فرزاد تهامی & استاد & \href{mailto:tahami@sharif.edu}{\lat{tahami@sharif.edu}} \\
\rowcolor{tableRow1}
دکتر مهران جاهد & دانشیار & \href{mailto:jahed@sharif.edu}{\lat{jahed@sharif.edu}} \\
\rowcolor{tableRow2}
دکتر خسرو حاج صادقی & دانشیار & \href{mailto:ksadeghi@sharif.edu}{\lat{ksadeghi@sharif.edu}} \\
\rowcolor{tableRow1}
دکتر احسان حاجی‌پور & استادیار & \href{mailto:e_hajpour@sharif.edu}{\lat{e\_hajpour@sharif.edu}} \\
\rowcolor{tableRow2}
دکتر صدیقه حاجی‌پور & استادیار & \href{mailto:hajlpour@sharif.edu}{\lat{hajlpour@sharif.edu}} \\
\rowcolor{tableRow1}
دکتر محمد حائری & استاد & \href{mailto:haeri@sharif.edu}{\lat{haeri@sharif.edu}} \\
\rowcolor{tableRow2}
دکتر بابک حسین خلج & استاد & \href{mailto:khalaj@sharif.edu}{\lat{khalaj@sharif.edu}} \\
\rowcolor{tableRow1}
دکتر سید حمید حسینی & استاد & \href{mailto:hosseini@sharif.edu}{\lat{hosseini@sharif.edu}} \\
\rowcolor{tableRow2}
دکتر یاس حسینی تهرانی & استادیار & \href{mailto:ytehrani@sharif.edu}{\lat{ytehrani@sharif.edu}} \\
\rowcolor{tableRow1}
دکتر امین خواسی & دانشیار & \href{mailto:khavasi@sharif.ir}{\lat{khavasi@sharif.ir}} \\
\rowcolor{tableRow2}
دکتر محمد امیر دستغیب & استادیار & \href{mailto:dastgheib@sharif.edu}{\lat{dastgheib@sharif.edu}} \\
\rowcolor{tableRow1}
دکتر محمدرضا زلقدری & دانشیار & \href{mailto:zolghadr@sharif.edu}{\lat{zolghadr@sharif.edu}} \\
\rowcolor{tableRow2}
دکتر محمدرضا روانجی & استادیار & \href{mailto:ravanji@sharif.edu}{\lat{ravanji@sharif.edu}} \\
\rowcolor{tableRow1}
دکتر بهزاد رجاعی سلامت & دانشیار & \href{mailto:rajaei@sharif.edu}{\lat{rajaei@sharif.edu}} \\
\rowcolor{tableRow2}
دکتر بیژن رشیدیان & استاد & \href{mailto:rashidian@sharif.edu}{\lat{rashidian@sharif.edu}} \\
\rowcolor{tableRow1}
دکتر امین رضایی‌زاده & استادیار & \href{mailto:aminre@sharif.ir}{\lat{aminre@sharif.ir}} \\
\rowcolor{tableRow2}
دکتر مصطفی زرغانی & استادیار & \href{mailto:zarghani@sharif.edu}{\lat{zarghani@sharif.edu}} \\
\rowcolor{tableRow1}
دکتر رضا سرواری & دانشیار & \href{mailto:sarvari@sharif.edu}{\lat{sarvari@sharif.edu}} \\
\rowcolor{tableRow2}
دکتر حامد شاه‌منصوری & استادیار & \href{mailto:hamedsh@sharif.edu}{\lat{hamedsh@sharif.edu}} \\
\rowcolor{tableRow1}
دکتر محمد شریف‌خانی & دانشیار & \href{mailto:msharifk@sharif.edu}{\lat{msharifk@sharif.edu}} \\
\rowcolor{tableRow2}
دکتر مهدی شعبانی & دانشیار & \href{mailto:mahdi@sharif.edu}{\lat{mahdi@sharif.edu}} \\
\rowcolor{tableRow1}
دکتر محمدباقر شمس‌اللهی & استاد & \href{mailto:mbshams@sharif.edu}{\lat{mbshams@sharif.edu}} \\
\rowcolor{tableRow2}
دکتر امیراحمد شیشه‌گر & استاد & \href{mailto:shishegar@sharif.edu}{\lat{shishegar@sharif.edu}} \\
\rowcolor{tableRow1}
دکتر جواد صالحی & استاد & \href{mailto:jsalehi@sharif.edu}{\lat{jsalehi@sharif.edu}} \\
\rowcolor{tableRow2}
دکتر امیر صفدریان & دانشیار & \href{mailto:safdarian@sharif.edu}{\lat{safdarian@sharif.edu}} \\
\rowcolor{tableRow1}
دکتر محمدرضا عارف & استاد & \href{mailto:aref@sharif.edu}{\lat{aref@sharif.edu}} \\
\rowcolor{tableRow2}
دکتر سید مجتبی عطارودی & استاد & \href{mailto:atarodi@sharif.edu}{\lat{atarodi@sharif.edu}} \\
\rowcolor{tableRow1}
دکتر عماد فاطمی‌زاده & دانشیار & \href{mailto:fatemizadeh@sharif.edu}{\lat{fatemizadeh@sharif.edu}} \\
\rowcolor{tableRow2}
دکتر علی فتوح احمدی & دانشیار & \href{mailto:afotowat@sharif.edu}{\lat{afotowat@sharif.edu}} \\
\rowcolor{tableRow1}
دکتر محمود فتوحی فیروزآبادی & استاد & \href{mailto:fotuhi@sharif.edu}{\lat{fotuhi@sharif.edu}} \\
\rowcolor{tableRow2}
دکتر محمد فخارزاده & دانشیار & \href{mailto:fakharzadeh@sharif.edu}{\lat{fakharzadeh@sharif.edu}} \\
\rowcolor{tableRow1}
دکتر مهدی فردمنش & استاد & \href{mailto:fardmanesh@sharif.edu}{\lat{fardmanesh@sharif.edu}} \\
\rowcolor{tableRow2}
دکتر فروهر فرزانه & استاد & \href{mailto:farzaneh@sharif.edu}{\lat{farzaneh@sharif.edu}} \\
\rowcolor{tableRow1}
دکتر علیرضا فرهادی & دانشیار & \href{mailto:afarhadi@sharif.edu}{\lat{afarhadi@sharif.edu}} \\
\rowcolor{tableRow2}
دکتر علی غضنفری & دانشیار & \href{mailto:ghazizadeh@sharif.edu}{\lat{ghazizadeh@sharif.edu}} \\
\rowcolor{tableRow1}
دکتر شهریار کابلی & استاد & \href{mailto:kaboli@sharif.edu}{\lat{kaboli@sharif.edu}} \\
\rowcolor{tableRow2}
دکتر زهرا کاوه‌وش & استاد & \href{mailto:kavehvash@sharif.edu}{\lat{kavehvash@sharif.edu}} \\
\rowcolor{tableRow1}
دکتر سید محمد کرباسی & دانشیار & \href{mailto:m.karbasi@sharif.ir}{\lat{m.karbasi@sharif.ir}} \\
\rowcolor{tableRow2}
دکتر سید جمال‌الدین گلستانی & دانشیار & \href{mailto:golestani@sharif.edu}{\lat{golestani@sharif.edu}} \\
\rowcolor{tableRow1}
دکتر محمد مهدی مجاهدین & استادیار & \href{mailto:mojahedian@sharif.edu}{\lat{mojahedian@sharif.edu}} \\
\rowcolor{tableRow2}
دکتر هدی محمدزاده & دانشیار & \href{mailto:hoda@sharif.edu}{\lat{hoda@sharif.edu}} \\
\rowcolor{tableRow1}
دکتر حسین مختاری & استاد & \href{mailto:mokhtari@sharif.edu}{\lat{mokhtari@sharif.edu}} \\
\rowcolor{tableRow2}
دکتر علی مدی & استاد & \href{mailto:medi@sharif.edu}{\lat{medi@sharif.edu}} \\
\rowcolor{tableRow1}
دکتر محمد ممریان & دانشیار & \href{mailto:mmemarian@sharif.ir}{\lat{mmemarian@sharif.ir}} \\
\rowcolor{tableRow2}
دکتر خاشع مهرانی & استاد & \href{mailto:mehrani@sharif.edu}{\lat{mehrani@sharif.edu}} \\
\rowcolor{tableRow1}
دکتر مهتاب میرمحسنی & دانشیار & \href{mailto:mirmohseni@sharif.edu}{\lat{mirmohseni@sharif.edu}} \\
\rowcolor{tableRow2}
دکتر مهرزاد نمور & دانشیار & \href{mailto:namvar@sharif.edu}{\lat{namvar@sharif.edu}} \\
\rowcolor{tableRow1}
دکتر محمد مهدی نایبی & استاد & \href{mailto:nayebi@sharif.edu}{\lat{nayebi@sharif.edu}} \\
\rowcolor{tableRow2}
دکتر زهرا نصیری قیداری & استاد & \href{mailto:znasiri@sharif.edu}{\lat{znasiri@sharif.edu}} \\
\rowcolor{tableRow1}
دکتر معصومه نصیری کناری & استاد & \href{mailto:mnasiri@sharif.edu}{\lat{mnasiri@sharif.edu}} \\
\rowcolor{tableRow2}
دکتر امین نوبختی & دانشیار & \href{mailto:nobakhti@sharif.edu}{\lat{nobakhti@sharif.edu}} \\
\rowcolor{tableRow1}
دکتر بیژن وثوقی وحدت & دانشیار & \href{mailto:vahdat@sharif.edu}{\lat{vahdat@sharif.edu}} \\
\rowcolor{tableRow2}
دکتر مهدی وکیلیان & استاد & \href{mailto:vakilian@sharif.edu}{\lat{vakilian@sharif.edu}} \\
\rowcolor{tableRow1}
دکتر محمد هادی & استادیار & \href{mailto:mohammad.hadi@sharif.edu}{\lat{mohammad.hadi@sharif.edu}} \\
\rowcolor{tableRow2}
دکتر متین هاشمی & دانشیار & \href{mailto:matin@sharif.edu}{\lat{matin@sharif.edu}} \\
\rowcolor{tableRow1}
دکتر محمدحسین یاسایی & استادیار & \href{mailto:yassaee@sharif.edu}{\lat{yassaee@sharif.edu}} \\
\bottomrule
\end{tabularx}
\begin{tablenotes}\small
\item این ایمیل‌ها از صفحه رسمی دانشکده مهندسی برق دانشگاه صنعتی شریف استخراج شده‌اند.
\end{tablenotes}
\end{threeparttable}
\end{table}

% ===================== CHAPTER 3: PHYSICS =====================
\clearpage
\chapter{اعضای هیئت علمی فیزیک}

\begin{keyconceptbox}
لیست کامل ایمیل‌های اعضای هیئت علمی دانشکده فیزیک دانشگاه صنعتی شریف. این اطلاعات از صفحه رسمی دانشکده استخراج شده‌اند.
\end{keyconceptbox}

\begin{table}[H]
\centering
\caption{لیست ایمیل‌های اعضای هیئت علمی فیزیک}
\label{tab:physics_emails}
\begin{threeparttable}
\begin{tabularx}{\textwidth}{H{4cm} X C{3.5cm}}
\toprule
\rowcolor{tableHeader}
\textcolor{white}{\textbf{نام}} & \textcolor{white}{\textbf{سمت}} & \textcolor{white}{\textbf{آدرس ایمیل}} \\
\midrule
\rowcolor{tableRow1}
دکتر علی ابوالحسنی & استادیار & \href{mailto:ali.abolhasani@sharif.edu}{\lat{ali.abolhasani@sharif.edu}} \\
\rowcolor{tableRow2}
دکتر امید اخوان & دانشیار & \href{mailto:oakhavan@sharif.edu}{\lat{oakhavan@sharif.edu}} \\
\rowcolor{tableRow1}
دکتر آرش عربی & استادیار & \href{mailto:a.a.ardehali@gmail.com}{\lat{a.a.ardehali@gmail.com}} \\
\rowcolor{tableRow2}
دکتر شانت باغرام & دانشیار & \href{mailto:baghram@sharif.edu}{\lat{baghram@sharif.edu}} \\
\rowcolor{tableRow1}
دکتر محمود بهمن‌آبادی & استاد & \href{mailto:bahmanabadi@sharif.edu}{\lat{bahmanabadi@sharif.edu}} \\
\rowcolor{tableRow2}
دکتر علیرضا بهرام‌پور & استاد & \href{mailto:bahrampour@sharif.edu}{\lat{bahrampour@sharif.edu}} \\
\rowcolor{tableRow1}
دکتر محمدرضا اجتهدی & استاد & \href{mailto:ejtehadi@sharif.edu}{\lat{ejtehadi@sharif.edu}} \\
\rowcolor{tableRow2}
دکتر علی اسفندیار & دانشیار & \href{mailto:esfandiar@physics.sharif.ir}{\lat{esfandiar@physics.sharif.ir}} \\
\rowcolor{tableRow1}
دکتر امین فرجی‌آستانه & استادیار & \href{mailto:faraji@sharif.ir}{\lat{faraji@sharif.ir}} \\
\rowcolor{tableRow2}
دکتر اعظم ایرجی‌زاد & استاد ممتاز & \href{mailto:iraji@sharif.edu}{\lat{iraji@sharif.edu}} \\
\rowcolor{tableRow1}
دکتر مهدی کارگیانی & دانشیار & \href{mailto:kargarian@sharif.edu}{\lat{kargarian@sharif.edu}} \\
\rowcolor{tableRow2}
دکتر وحید کریمی‌پور & استاد & \href{mailto:vahid@sharif.edu}{\lat{vahid@sharif.edu}} \\
\rowcolor{tableRow1}
دکتر علی خادمی & استادیار & \href{mailto:khademi@sharif.edu}{\lat{khademi@sharif.edu}} \\
\rowcolor{tableRow2}
دکتر نیما خسروی & دانشیار & \href{mailto:nima@sharif.edu}{\lat{nima@sharif.edu}} \\
\rowcolor{tableRow1}
دکتر عبدالله لنگری & استاد & \href{mailto:langari@sharif.edu}{\lat{langari@sharif.edu}} \\
\rowcolor{tableRow2}
دکتر لاله ممارزاده & دانشیار & \href{mailto:memarzadeh@sharif.edu}{\lat{memarzadeh@sharif.edu}} \\
\rowcolor{tableRow1}
دکتر مریم میرکمالی & استادیار & \href{mailto:mirkamali@sharif.edu}{\lat{mirkamali@sharif.edu}} \\
\rowcolor{tableRow2}
دکتر سامان مقیمی آراگی & دانشیار & \href{mailto:samanimi@sharif.edu}{\lat{samanimi@sharif.edu}} \\
\rowcolor{tableRow1}
دکتر علیرضا مشفق & استاد ممتاز & \href{mailto:moshfegh@sharif.edu}{\lat{moshfegh@sharif.edu}} \\
\rowcolor{tableRow2}
دکتر نعیمه ناصری & دانشیار & \href{mailto:naseri@sharif.edu}{\lat{naseri@sharif.edu}} \\
\rowcolor{tableRow1}
دکتر صادق رئیسی & استادیار & \href{mailto:sraeisi@sharif.edu}{\lat{sraeisi@sharif.edu}} \\
\rowcolor{tableRow2}
دکتر محمدرضا رحیمی‌تبار & استاد & \href{mailto:rahimitabar@sharif.edu}{\lat{rahimitabar@sharif.edu}} \\
\rowcolor{tableRow1}
دکتر سهراب رهبر & استاد & \href{mailto:rahvar@sharif.edu}{\lat{rahvar@sharif.edu}} \\
\rowcolor{tableRow2}
دکتر رضا رضایی & استادیار & \href{mailto:reza.rezaei@sharif.edu}{\lat{reza.rezaei@sharif.edu}} \\
\rowcolor{tableRow1}
دکتر علی رزاخانی & دانشیار & \href{mailto:rezakhani@sharif.edu}{\lat{rezakhani@sharif.edu}} \\
\rowcolor{tableRow2}
دکتر شاهین روحانی & استاد & \href{mailto:srouhani@sharif.edu}{\lat{srouhani@sharif.edu}} \\
\rowcolor{tableRow1}
دکتر ندا سدوقی & استاد & \href{mailto:sadooghi@sharif.ir}{\lat{sadooghi@sharif.ir}} \\
\rowcolor{tableRow2}
دکتر سید نادر سید ریحانی & دانشیار & \href{mailto:sreihani@sharif.ir}{\lat{sreihani@sharif.ir}} \\
\rowcolor{tableRow1}
دکتر نیما تقوی‌نیا & استاد & \href{mailto:taghavinia@sharif.edu}{\lat{taghavinia@sharif.edu}} \\
\rowcolor{tableRow2}
دکتر مهدی ترابیان & استادیار & \href{mailto:mahdi@physics.sharif.ir}{\lat{mahdi@physics.sharif.ir}} \\
\rowcolor{tableRow1}
دکتر سید می‌ر ابوالحسن واعظی & استادیار & \href{mailto:vaezi@sharif.edu}{\lat{vaezi@sharif.edu}} \\
\bottomrule
\end{tabularx}
\begin{tablenotes}\small
\item این ایمیل‌ها از صفحه رسمی دانشکده فیزیک دانشگاه صنعتی شریف استخراج شده‌اند.
\end{tablenotes}
\end{threeparttable}
\end{table}

% ===================== CHAPTER 4: COMPLETE EMAIL LIST (4 TABLES) =====================
\clearpage
\chapter{لیست کامل ایمیل‌ها (چهار بخش)}

\begin{keyconceptbox}
لیست تک‌ستونی تمامی ایمیل‌های اعضای هیئت علمی دانشکده‌های مهندسی برق و فیزیک. این لیست به چهار جدول تقسیم شده برای خوانایی بهتر و جلوگیری از سرریز صفحات. مناسب برای کپی و استفاده در نرم‌افزارهای ارسال ایمیل گروهی.
\end{keyconceptbox}

\section{جدول ۱: ایمیل‌های ۱ تا ۲۶}

\begin{table}[H]
\centering
\caption{لیست تک‌ستونی ایمیل‌ها — بخش اول}
\label{tab:complete_email_list_1}
\begin{threeparttable}
\begin{tabularx}{\textwidth}{>{\raggedright}X}
\toprule
\rowcolor{tableHeader}
\textcolor{white}{\textbf{آدرس‌های ایمیل (قابل کپی)}} \\
\midrule
\rowcolor{tableRow1}
\href{mailto:ashtiani@sharif.edu}{\lat{ashtiani@sharif.edu}} \\
\rowcolor{tableRow2}
\href{mailto:ahi@sharif.edu}{\lat{ahi@sharif.edu}} \\
\rowcolor{tableRow1}
\href{mailto:ahmadi@sharif.edu}{\lat{ahmadi@sharif.edu}} \\
\rowcolor{tableRow2}
\href{mailto:fakbar@sharif.edu}{\lat{fakbar@sharif.edu}} \\
\rowcolor{tableRow1}
\href{mailto:makbari@sharif.edu}{\lat{makbari@sharif.edu}} \\
\rowcolor{tableRow2}
\href{mailto:amirl@sharif.edu}{\lat{amirl@sharif.edu}} \\
\rowcolor{tableRow1}
\href{mailto:aamini@sharif.edu}{\lat{aamini@sharif.edu}} \\
\rowcolor{tableRow2}
\href{mailto:oraee@sharif.edu}{\lat{oraee@sharif.edu}} \\
\rowcolor{tableRow1}
\href{mailto:babazadeh@sharif.ir}{\lat{babazadeh@sharif.ir}} \\
\rowcolor{tableRow2}
\href{mailto:mbzadeh@sharif.edu}{\lat{mbzadeh@sharif.edu}} \\
\rowcolor{tableRow1}
\href{mailto:bagheri-s@sharif.edu}{\lat{bagheri-s@sharif.edu}} \\
\rowcolor{tableRow2}
\href{mailto:banai@sharif.edu}{\lat{banai@sharif.edu}} \\
\rowcolor{tableRow1}
\href{mailto:behroozi@sharif.edu}{\lat{behroozi@sharif.edu}} \\
\rowcolor{tableRow2}
\href{mailto:behnia@sharif.edu}{\lat{behnia@sharif.edu}} \\
\rowcolor{tableRow1}
\href{mailto:pakravan@sharif.edu}{\lat{pakravan@sharif.edu}} \\
\rowcolor{tableRow2}
\href{mailto:parniani@sharif.edu}{\lat{parniani@sharif.edu}} \\
\rowcolor{tableRow1}
\href{mailto:tavazoei@sharif.edu}{\lat{tavazoei@sharif.edu}} \\
\rowcolor{tableRow2}
\href{mailto:tahami@sharif.edu}{\lat{tahami@sharif.edu}} \\
\rowcolor{tableRow1}
\href{mailto:jahed@sharif.edu}{\lat{jahed@sharif.edu}} \\
\rowcolor{tableRow2}
\href{mailto:ksadeghi@sharif.edu}{\lat{ksadeghi@sharif.edu}} \\
\rowcolor{tableRow1}
\href{mailto:e_hajpour@sharif.edu}{\lat{e\_hajpour@sharif.edu}} \\
\rowcolor{tableRow2}
\href{mailto:hajlpour@sharif.edu}{\lat{hajlpour@sharif.edu}} \\
\rowcolor{tableRow1}
\href{mailto:haeri@sharif.edu}{\lat{haeri@sharif.edu}} \\
\rowcolor{tableRow2}
\href{mailto:khalaj@sharif.edu}{\lat{khalaj@sharif.edu}} \\
\rowcolor{tableRow1}
\href{mailto:hosseini@sharif.edu}{\lat{hosseini@sharif.edu}} \\
\rowcolor{tableRow2}
\href{mailto:ytehrani@sharif.edu}{\lat{ytehrani@sharif.edu}} \\
\bottomrule
\end{tabularx}
\begin{tablenotes}\small
\item ایمیل‌های ۱ تا ۲۶ از لیست کامل ۱۰۴ تایی
\end{tablenotes}
\end{threeparttable}
\end{table}

\clearpage
\section{جدول ۲: ایمیل‌های ۲۷ تا ۵۲}

\begin{table}[H]
\centering
\caption{لیست تک‌ستونی ایمیل‌ها — بخش دوم}
\label{tab:complete_email_list_2}
\begin{threeparttable}
\begin{tabularx}{\textwidth}{>{\raggedright}X}
\toprule
\rowcolor{tableHeader}
\textcolor{white}{\textbf{آدرس‌های ایمیل (قابل کپی)}} \\
\midrule
\rowcolor{tableRow1}
\href{mailto:khavasi@sharif.ir}{\lat{khavasi@sharif.ir}} \\
\rowcolor{tableRow2}
\href{mailto:dastgheib@sharif.edu}{\lat{dastgheib@sharif.edu}} \\
\rowcolor{tableRow1}
\href{mailto:zolghadr@sharif.edu}{\lat{zolghadr@sharif.edu}} \\
\rowcolor{tableRow2}
\href{mailto:ravanji@sharif.edu}{\lat{ravanji@sharif.edu}} \\
\rowcolor{tableRow1}
\href{mailto:rajaei@sharif.edu}{\lat{rajaei@sharif.edu}} \\
\rowcolor{tableRow2}
\href{mailto:rashidian@sharif.edu}{\lat{rashidian@sharif.edu}} \\
\rowcolor{tableRow1}
\href{mailto:aminre@sharif.ir}{\lat{aminre@sharif.ir}} \\
\rowcolor{tableRow2}
\href{mailto:zarghani@sharif.edu}{\lat{zarghani@sharif.edu}} \\
\rowcolor{tableRow1}
\href{mailto:sarvari@sharif.edu}{\lat{sarvari@sharif.edu}} \\
\rowcolor{tableRow2}
\href{mailto:hamedsh@sharif.edu}{\lat{hamedsh@sharif.edu}} \\
\rowcolor{tableRow1}
\href{mailto:msharifk@sharif.edu}{\lat{msharifk@sharif.edu}} \\
\rowcolor{tableRow2}
\href{mailto:mahdi@sharif.edu}{\lat{mahdi@sharif.edu}} \\
\rowcolor{tableRow1}
\href{mailto:mbshams@sharif.edu}{\lat{mbshams@sharif.edu}} \\
\rowcolor{tableRow2}
\href{mailto:shishegar@sharif.edu}{\lat{shishegar@sharif.edu}} \\
\rowcolor{tableRow1}
\href{mailto:jsalehi@sharif.edu}{\lat{jsalehi@sharif.edu}} \\
\rowcolor{tableRow2}
\href{mailto:safdarian@sharif.edu}{\lat{safdarian@sharif.edu}} \\
\rowcolor{tableRow1}
\href{mailto:aref@sharif.edu}{\lat{aref@sharif.edu}} \\
\rowcolor{tableRow2}
\href{mailto:atarodi@sharif.edu}{\lat{atarodi@sharif.edu}} \\
\rowcolor{tableRow1}
\href{mailto:fatemizadeh@sharif.edu}{\lat{fatemizadeh@sharif.edu}} \\
\rowcolor{tableRow2}
\href{mailto:afotowat@sharif.edu}{\lat{afotowat@sharif.edu}} \\
\rowcolor{tableRow1}
\href{mailto:fotuhi@sharif.edu}{\lat{fotuhi@sharif.edu}} \\
\rowcolor{tableRow2}
\href{mailto:fakharzadeh@sharif.edu}{\lat{fakharzadeh@sharif.edu}} \\
\rowcolor{tableRow1}
\href{mailto:fardmanesh@sharif.edu}{\lat{fardmanesh@sharif.edu}} \\
\rowcolor{tableRow2}
\href{mailto:farzaneh@sharif.edu}{\lat{farzaneh@sharif.edu}} \\
\rowcolor{tableRow1}
\href{mailto:afarhadi@sharif.edu}{\lat{afarhadi@sharif.edu}} \\
\rowcolor{tableRow2}
\href{mailto:ghazizadeh@sharif.edu}{\lat{ghazizadeh@sharif.edu}} \\
\bottomrule
\end{tabularx}
\begin{tablenotes}\small
\item ایمیل‌های ۲۷ تا ۵۲ از لیست کامل ۱۰۴ تایی
\end{tablenotes}
\end{threeparttable}
\end{table}

\clearpage
\section{جدول ۳: ایمیل‌های ۵۳ تا ۷۸}

\begin{table}[H]
\centering
\caption{لیست تک‌ستونی ایمیل‌ها — بخش سوم}
\label{tab:complete_email_list_3}
\begin{threeparttable}
\begin{tabularx}{\textwidth}{>{\raggedright}X}
\toprule
\rowcolor{tableHeader}
\textcolor{white}{\textbf{آدرس‌های ایمیل (قابل کپی)}} \\
\midrule
\rowcolor{tableRow1}
\href{mailto:kaboli@sharif.edu}{\lat{kaboli@sharif.edu}} \\
\rowcolor{tableRow2}
\href{mailto:kavehvash@sharif.edu}{\lat{kavehvash@sharif.edu}} \\
\rowcolor{tableRow1}
\href{mailto:m.karbasi@sharif.ir}{\lat{m.karbasi@sharif.ir}} \\
\rowcolor{tableRow2}
\href{mailto:golestani@sharif.edu}{\lat{golestani@sharif.edu}} \\
\rowcolor{tableRow1}
\href{mailto:mojahedian@sharif.edu}{\lat{mojahedian@sharif.edu}} \\
\rowcolor{tableRow2}
\href{mailto:hoda@sharif.edu}{\lat{hoda@sharif.edu}} \\
\rowcolor{tableRow1}
\href{mailto:mokhtari@sharif.edu}{\lat{mokhtari@sharif.edu}} \\
\rowcolor{tableRow2}
\href{mailto:medi@sharif.edu}{\lat{medi@sharif.edu}} \\
\rowcolor{tableRow1}
\href{mailto:mmemarian@sharif.ir}{\lat{mmemarian@sharif.ir}} \\
\rowcolor{tableRow2}
\href{mailto:mehrani@sharif.edu}{\lat{mehrani@sharif.edu}} \\
\rowcolor{tableRow1}
\href{mailto:mirmohseni@sharif.edu}{\lat{mirmohseni@sharif.edu}} \\
\rowcolor{tableRow2}
\href{mailto:namvar@sharif.edu}{\lat{namvar@sharif.edu}} \\
\rowcolor{tableRow1}
\href{mailto:nayebi@sharif.edu}{\lat{nayebi@sharif.edu}} \\
\rowcolor{tableRow2}
\href{mailto:znasiri@sharif.edu}{\lat{znasiri@sharif.edu}} \\
\rowcolor{tableRow1}
\href{mailto:mnasiri@sharif.edu}{\lat{mnasiri@sharif.edu}} \\
\rowcolor{tableRow2}
\href{mailto:nobakhti@sharif.edu}{\lat{nobakhti@sharif.edu}} \\
\rowcolor{tableRow1}
\href{mailto:vahdat@sharif.edu}{\lat{vahdat@sharif.edu}} \\
\rowcolor{tableRow2}
\href{mailto:vakilian@sharif.edu}{\lat{vakilian@sharif.edu}} \\
\rowcolor{tableRow1}
\href{mailto:mohammad.hadi@sharif.edu}{\lat{mohammad.hadi@sharif.edu}} \\
\rowcolor{tableRow2}
\href{mailto:matin@sharif.edu}{\lat{matin@sharif.edu}} \\
\rowcolor{tableRow1}
\href{mailto:yassaee@sharif.edu}{\lat{yassaee@sharif.edu}} \\
\rowcolor{tableRow2}
\href{mailto:ali.abolhasani@sharif.edu}{\lat{ali.abolhasani@sharif.edu}} \\
\rowcolor{tableRow1}
\href{mailto:oakhavan@sharif.edu}{\lat{oakhavan@sharif.edu}} \\
\rowcolor{tableRow2}
\href{mailto:a.a.ardehali@gmail.com}{\lat{a.a.ardehali@gmail.com}} \\
\rowcolor{tableRow1}
\href{mailto:baghram@sharif.edu}{\lat{baghram@sharif.edu}} \\
\rowcolor{tableRow2}
\href{mailto:bahmanabadi@sharif.edu}{\lat{bahmanabadi@sharif.edu}} \\
\bottomrule
\end{tabularx}
\begin{tablenotes}\small
\item ایمیل‌های ۵۳ تا ۷۸ از لیست کامل ۱۰۴ تایی
\end{tablenotes}
\end{threeparttable}
\end{table}

\clearpage
\section{جدول ۴: ایمیل‌های ۷۹ تا ۱۰۴}

\begin{table}[H]
\centering
\caption{لیست تک‌ستونی ایمیل‌ها — بخش پایانی}
\label{tab:complete_email_list_4}
\begin{threeparttable}
\begin{tabularx}{\textwidth}{>{\raggedright}X}
\toprule
\rowcolor{tableHeader}
\textcolor{white}{\textbf{آدرس‌های ایمیل (قابل کپی)}} \\
\midrule
\rowcolor{tableRow1}
\href{mailto:bahrampour@sharif.edu}{\lat{bahrampour@sharif.edu}} \\
\rowcolor{tableRow2}
\href{mailto:ejtehadi@sharif.edu}{\lat{ejtehadi@sharif.edu}} \\
\rowcolor{tableRow1}
\href{mailto:esfandiar@physics.sharif.ir}{\lat{esfandiar@physics.sharif.ir}} \\
\rowcolor{tableRow2}
\href{mailto:faraji@sharif.ir}{\lat{faraji@sharif.ir}} \\
\rowcolor{tableRow1}
\href{mailto:iraji@sharif.edu}{\lat{iraji@sharif.edu}} \\
\rowcolor{tableRow2}
\href{mailto:kargarian@sharif.edu}{\lat{kargarian@sharif.edu}} \\
\rowcolor{tableRow1}
\href{mailto:vahid@sharif.edu}{\lat{vahid@sharif.edu}} \\
\rowcolor{tableRow2}
\href{mailto:khademi@sharif.edu}{\lat{khademi@sharif.edu}} \\
\rowcolor{tableRow1}
\href{mailto:nima@sharif.edu}{\lat{nima@sharif.edu}} \\
\rowcolor{tableRow2}
\href{mailto:langari@sharif.edu}{\lat{langari@sharif.edu}} \\
\rowcolor{tableRow1}
\href{mailto:memarzadeh@sharif.edu}{\lat{memarzadeh@sharif.edu}} \\
\rowcolor{tableRow2}
\href{mailto:mirkamali@sharif.edu}{\lat{mirkamali@sharif.edu}} \\
\rowcolor{tableRow1}
\href{mailto:samanimi@sharif.edu}{\lat{samanimi@sharif.edu}} \\
\rowcolor{tableRow2}
\href{mailto:moshfegh@sharif.edu}{\lat{moshfegh@sharif.edu}} \\
\rowcolor{tableRow1}
\href{mailto:naseri@sharif.edu}{\lat{naseri@sharif.edu}} \\
\rowcolor{tableRow2}
\href{mailto:sraeisi@sharif.edu}{\lat{sraeisi@sharif.edu}} \\
\rowcolor{tableRow1}
\href{mailto:rahimitabar@sharif.edu}{\lat{rahimitabar@sharif.edu}} \\
\rowcolor{tableRow2}
\href{mailto:rahvar@sharif.edu}{\lat{rahvar@sharif.edu}} \\
\rowcolor{tableRow1}
\href{mailto:reza.rezaei@sharif.edu}{\lat{reza.rezaei@sharif.edu}} \\
\rowcolor{tableRow2}
\href{mailto:rezakhani@sharif.edu}{\lat{rezakhani@sharif.edu}} \\
\rowcolor{tableRow1}
\href{mailto:srouhani@sharif.edu}{\lat{srouhani@sharif.edu}} \\
\rowcolor{tableRow2}
\href{mailto:sadooghi@sharif.ir}{\lat{sadooghi@sharif.ir}} \\
\rowcolor{tableRow1}
\href{mailto:sreihani@sharif.ir}{\lat{sreihani@sharif.ir}} \\
\rowcolor{tableRow2}
\href{mailto:taghavinia@sharif.edu}{\lat{taghavinia@sharif.edu}} \\
\rowcolor{tableRow1}
\href{mailto:mahdi@physics.sharif.ir}{\lat{mahdi@physics.sharif.ir}} \\
\rowcolor{tableRow2}
\href{mailto:vaezi@sharif.edu}{\lat{vaezi@sharif.edu}} \\
\bottomrule
\end{tabularx}
\begin{tablenotes}\small
\item ایمیل‌های ۷۹ تا ۱۰۴ (آخرین ایمیل) از لیست کامل ۱۰۴ تایی
\end{tablenotes}
\end{threeparttable}
\end{table}

% ===================== FINAL PAGE =====================
\clearpage
\thispagestyle{empty}
\begin{center}
\vspace*{3cm}

{\Huge\bfseries\color{headerColor}پایان سند}

\vspace{1cm}
{\color{headerColor}\rule{0.6\textwidth}{0.75pt}}

\vspace{1cm}
{\Large\bfseries لیست تماس اعضای هیئت علمی}

\vspace{0.5cm}
{\large دانشگاه صنعتی شریف}

\vspace{2cm}
{\color{gray} به‌روزرسانی: بهمن ۱۴۰۴}
\end{center}

\end{document}